\SourceLeaf{173}\bigskip
\Needspace{10\baselineskip}
\centerline{[3.]}

\smallskip

\centerline{\textls{Ideas.}}

\smallskip

The most precise determination of stellar declinations could be made near the equator with good repeating theodolites by observing the differences in azimuth between the stars and terrestrial objects at the times of the stars' greatest eastern and western elongations. \hfill\mbox{1813 April 21.}

\smallskip

In the theory of parallel lines we have yet to advance beyond Euclid. This is the shameful part of mathematics and must sooner or later assume an entirely different form. \hfill\mbox{1813 April 27.}

\smallskip

Just as a class of truths of higher arithmetic, restricted to real numbers, is closely connected with circular functions, so the more general truths of H.\,A., extended to imaginary quantities, will as it were be reflected in the lemniscatic functions. These ideas open the gate to an exceedingly rich field of analysis.

\smallskip

The crystal prism on telescopes increases their length by about $\frac{1}{3}$
 of the side of the prism, but somewhat impairs their perfect colorlessness.

\bigskip
\hrule
\bigskip

\centerline{REMARK.}

\smallskip

\textsc{Schumacher} stayed in Göttingen during the winter of 1808 to 1809. He recorded \textsc{Gauss}'s conversations with him in a notebook entitled »GAUSSIANA«. Note [2.] is taken from this notebook, while note [3.] is written on a separate slip of paper.
\begin{flushright}\textsc{Stäckel}.\end{flushright}

\bigskip
\hrule
